% 5. สรุปผลและข้อเสนอแนะ -- ย่อจากรายงานบทที่ 5 (../report/chapter/05-conclusion.tex)
% ตัวเลข power คำนวณจากการแปลง Fisher z ที่ alpha = 0.05 สองทาง (เหมือนในรายงาน)
\section{สรุปผลและข้อเสนอแนะ}

\subsection{สรุปผล}

ในนักศึกษา 107 คน ชั่วโมงใช้โซเชียลมีเดียต่อวันไม่สัมพันธ์กับภาวะซึมเศร้า ความวิตกกังวล
และความเครียด ทั้งก่อนและหลังควบคุมปัจจัยอื่น ประเภทเนื้อหาที่ดูมากที่สุดไม่ทำให้คะแนนต่างกัน
และการติดตามข่าวเชิงลบสัมพันธ์กับความวิตกกังวลมากกว่ามิติอื่นตามที่คาดไว้ แต่ไม่มีนัยสำคัญ
ตัวแปรที่สัมพันธ์กับคะแนนชัดที่สุดคือความเพียงพอทางการเงิน
ตามด้วยคุณภาพการนอนและความกดดันจากการเรียน

\subsection{อภิปรายผล}

ผลที่ไม่พบความสัมพันธ์กับเวลาที่ใช้สอดคล้องกับงานของ Orben และ Przybylski
\cite{orben2019association} ซึ่งพบขนาดผลเพียงร้อยละ 0.4 หรือ \stat{r} ประมาณ 0.06
หากความสัมพันธ์จริงมีขนาดเท่านี้ กลุ่มตัวอย่าง 107 คนมีโอกาสตรวจพบเพียงประมาณร้อยละ 10
และต้องใช้ผู้ตอบราว 1,960 คนจึงจะมีโอกาสตรวจพบร้อยละ 80
ผลนี้จึงไม่ได้แปลว่าไม่มีความสัมพันธ์ แต่บอกว่าถ้ามีก็เล็กเกินกว่ากลุ่มตัวอย่างขนาดนี้จะตรวจพบ
นอกจากนี้ ผู้ตอบทุกคนใช้โซเชียลมีเดียอย่างน้อย 4 ชั่วโมงต่อวัน
ช่วงของตัวแปรที่แคบเช่นนี้ทำให้ค่า \stat{r} เล็กลงได้
ผลนี้ยังเข้ากับข้อสรุปของ Sala และคณะ \cite{sala2024umbrella}
ว่าผลไม่ได้ขึ้นกับเวลาที่ใช้อย่างเดียว

ทิศทางของผลเรื่องการติดตามข่าวตรงกับงานของ Dominguez-Rodriguez และคณะ
\cite{dominguez2025climate} แต่ไม่มีนัยสำคัญ เหตุผลที่เป็นไปได้มีสามข้อ
ข้อแรก งานวิจัยนี้วัดเพียงความถี่ของการติดตามข่าว ไม่ได้ใช้แบบวัด doomscrolling เต็มรูปแบบ
ข้อที่สอง มิติความวิตกกังวลมีความเชื่อมั่นต่ำที่สุด ซึ่งทำให้ค่า \stat{r} ที่วัดได้ต่ำกว่าความจริง
ข้อที่สาม ถ้าขนาดผลจริงเท่ากับที่พบ (\stat{r}~= 0.146)
ต้องใช้ผู้ตอบราว 366 คนจึงจะมีโอกาสตรวจพบร้อยละ 80

ความเพียงพอทางการเงินถูกใส่ไว้ในฐานะตัวแปรควบคุม
แต่กลับเป็นตัวแปรที่แยกคะแนนภาวะซึมเศร้าได้ชัดที่สุด
อย่างไรก็ตาม ข้อมูลเก็บครั้งเดียว จึงบอกไม่ได้ว่าปัญหาการเงินทำให้ซึมเศร้า
หรือผู้ที่ซึมเศร้าประเมินสภาพการเงินของตนแย่กว่าความจริง
ส่วนกลไกการแย่งเวลาไม่ได้รับการสนับสนุน เพราะชั่วโมงใช้งานไม่สัมพันธ์กับการนอน
และสิ่งที่สัมพันธ์กับคะแนนคือคุณภาพการนอน ไม่ใช่จำนวนชั่วโมง
ด้านแบบวัด ค่าแอลฟาทั้งฉบับ (0.894) ใกล้กับที่ Manzar และคณะ \cite{manzar2025dass21} รายงานไว้
แต่ภาวะซึมเศร้ากับความเครียดสัมพันธ์กันค่อนข้างสูง (\stat{r}~= 0.68)
การเทียบว่าตัวแปรใดสัมพันธ์กับมิติใดมากกว่ากันจึงต้องตีความอย่างระมัดระวัง

\subsection{ข้อจำกัด}

ข้อมูลเก็บครั้งเดียว จึงบอกได้เพียงความสัมพันธ์ ไม่ใช่เหตุและผล
กลุ่มตัวอย่างเลือกแบบตามสะดวกจากมหาวิทยาลัยเดียว มีจำนวน 107 คน น้อยกว่าเป้า 150 ถึง 200 คน
และกระจุกในชั้นปีที่ 2 (ร้อยละ 81.3) กลุ่มเนื้อหาข่าวมีเพียง 8 คน
ตัวแปรหลายตัวเป็นการรายงานตนเอง และการปรับค่า \stat{p} ด้วยวิธี Bonferroni ทำเฉพาะการถดถอย
ผลการศึกษาใช้ได้ในระดับกลุ่มเท่านั้น ไม่ใช้วินิจฉัยสุขภาพจิตรายบุคคล

\subsection{ข้อเสนอแนะ}

ในการนำผลไปใช้ ข้อมูลชุดนี้สนับสนุนให้การดูแลสุขภาพจิตนักศึกษา
ให้ความสำคัญกับเรื่องค่าใช้จ่ายและคุณภาพการนอนมากกว่าการจำกัดเวลาหน้าจอ
สำหรับการศึกษาครั้งต่อไป ควร
(1)~กำหนดขนาดกลุ่มตัวอย่างจากขนาดผลที่คาดไว้ เช่น อย่างน้อย 366 คนสำหรับเรื่องการติดตามข่าว
(2)~สุ่มตัวอย่างให้กระจายตามชั้นปีและคณะ
(3)~ใช้แบบวัด doomscrolling เต็มรูปแบบ และวัดรูปแบบการใช้งานให้ละเอียดขึ้น
และ (4)~เก็บข้อมูลกลุ่มเดิมซ้ำหลายช่วงเวลา เพื่อดูว่าพฤติกรรมหรืออาการเกิดก่อน
